\section{Background: the VOUCH certificate}\label{sec:background}

This section recalls the parts of VOUCH \citep{dang2026vouch} that the extension builds
on and fixes notation; proofs are in the original paper.

\paragraph{Pipeline and threat model.} A provider fine-tunes a model on a corpus,
receives deletion requests, applies an unlearning routine to a forget set, and deploys
the result $M$. An auditor (the verifier) wants evidence that the routine removed the
influence of the forget set. The provider is \emph{honest but verifiable}: it runs the
committed pipeline on the committed data, and in particular processes the canaries
exactly as it processes the rest of the forget set, but it is not trusted to report
whether unlearning succeeded.

\paragraph{Paired ghost canaries.} Before fine-tuning, the verifier generates a cohort of
$m$ pairs with contents $C_i=(c_i^0,c_i^1)$ and draws independent fair coins
$b_i\sim\Bern(\frac12)$. The in-twin $c_i^{b_i}$ is inserted into the training corpus
and later into the forget set; the ghost twin $c_i^{1-b_i}$ is never trained on. The
manifest of contents and coins is committed (VOUCH uses a Merkle commitment), and each
coin is opened only when its pair is audited. All remaining randomness of training and
unlearning is collected in $\xi$.

\paragraph{Scores, signs and the certified quantity.} A score $s(M,c)$ measures how much
the model favours a record, for example the mean token log-probability of a secret
after its prompt. A declared class $\F$ of scores is fixed before any query. For each
pair and score the verifier computes the difference
$D_i^{(s)}=s(M,c_i^{\mathrm{in}})-s(M,c_i^{\mathrm{ghost}})$ and the sign
$Z_i^{(s)}=\mathbf 1[D_i^{(s)}>0]$, ties split by a committed coin. The certified
quantity is the realised cohort advantage
\[
\Delta^{(s)} \;=\; \frac{2}{m}\sum_{i=1}^m Z_i^{(s)} \;-\; 1 ,
\]
which is zero in expectation under exact unlearning and positive when the in-twins are
still favoured. Negative values (the in-twin scores conspicuously below its partner)
are leakage too, and VOUCH certifies both directions.

\paragraph{The exact null.} If $M$ is conditionally independent of the coins given the
contents, which holds for exact unlearning and for retraining from scratch, then for
every score and every step $t$ of a reveal order chosen independently of the signs, the
sign of the pair opened at step $t$ is a fair coin given everything opened before
(VOUCH Theorem~1). No independence across pairs is needed: the signs share one model,
but the coin governing each pair is fair given the past.

\paragraph{The certificate.} The verifier reveals the cohort in a uniformly random order
$\pi$, drawn independently of the signs (VOUCH uses a public randomness beacon), and
for each score and direction runs a betting wealth process
$W_t=\prod_{u\le t}\bigl(1+\lambda_u(p_0(u)-Z_{\pi(u)})\bigr)$ against the
\emph{without-replacement boundary}
\begin{equation}\label{eq:worboundary}
p_0(t) \;=\; \frac{m p_0 - S_{t-1}}{m-t+1},\qquad p_0=\tfrac12+\tfrac\eps2,
\end{equation}
where $S_{t-1}$ is the number of ones revealed before step $t$. If the cohort
advantage were at least $\eps$, the conditional mean of the next sign would be at least
$p_0(t)$, so $W_t$ is a nonnegative supermartingale and Ville's inequality bounds the
probability that it ever reaches $1/\alpha$. A two-sided $(\eps,\alpha)$-certificate is
issued when every per-score, per-direction process has crossed $1/\alpha$, and then
$\Pr[\text{a false certificate is ever issued}]\le\alpha$ under arbitrary dependence
between pairs (VOUCH Theorem~2). If $p_0(t)>1$ the null is infeasible and is refuted
outright. A separate revocation process tests the exact null directly and raises an
alarm (status REVOKED) when it accumulates evidence that the model still depends on the
coins. A run that neither certifies nor revokes ends UNDETERMINED.

\paragraph{Transfer to the coin-averaged advantage.} The cohort advantage depends on the
realised coins. If flipping one coin changes at most $\kappa$ signs, McDiarmid's
inequality \citep{mcdiarmid1989bounded} bounds the distance to the advantage averaged
over the coins with the contents held fixed by $\kappa\sqrt{2\log(2/\alpha')/m}$ (VOUCH
Proposition~3). Section~\ref{sec:organic} extends this transfer to the contents and then
to organic records.

\paragraph{Notation for this paper.} We write $\Delta$ for the cohort advantage of a
fixed score when the score is clear, $\F_t$ for the verifier's filtration after $t$
reveals, and $\mathbb{E}$ for expectation. ``Literal'' refers to scores computed on
the trained surface form of a canary, ``paraphrase'' to scores computed on questions
never seen in training.
